\section{Introduction}\label{sec:introduction}
A numerical policy and an error bound for that policy are different economic objects. An approximation to a value function can improve while the associated consumption, investment, or adjustment decision remains unchanged. This distinction matters for neural dynamic programming: a training criterion that reduces an upper bound on regret need not increase the payoff of the delivered policy. A useful numerical method must identify both changes rather than treating either one as evidence for the other.

We develop Neural Bellman Operators with value-separated policy acceptance. A neural network proposes controls. A policy evaluator supplies a lower and an upper bound for the payoff of those controls under the original stopping contract. An optimal-comparison witness supplies a separate upper bound on the value of the unrestricted economic problem. A policy update is accepted only when the new payoff lower bound exceeds the incumbent payoff upper bound on the declared comparison set. Updating the optimal-comparison witness does not, by itself, change the incumbent policy. The resulting distinction is operational: the acceptance test concerns an actual policy-value difference, whereas the optimality test concerns the distance from that policy to the original optimum.

Our main new analytical result makes policy-sensitive comparison useful beyond a single initial condition. For budget-feasible time controls in the original stochastic economy, a coupled-payoff calculation bounds the variation of a policy-value difference across an initial-state rectangle. It uses differences between the two policies, rather than adding separate, much larger value-function Lipschitz constants. Original CRRA utility, common Gaussian shocks, exact terminal-wealth accounting, and explicit stopping corrections enter the calculation. An independently enclosed value difference at the reference state therefore yields a uniform policy-improvement certificate on a continuum of initial states. A related, observable progress inequality certifies contraction of true policy regret without substituting a critic loss for that regret.

The numerical implementation trains new neural control generators from model-only random starts. No inherited optimal policy, dual witness, or solution label enters their training objective. Five seeds, four prescribed work budgets, and the independent payoff acceptance rule are fixed before execution. All fifteen retained transitions increase certified policy values at the reference state. Initial-to-final improvement is uniform on $K=[1.98,2.02]\times[1.24,1.26]$: the gain is at least $\NineteenGain$ utility units for every reported seed. Combining this gain with the independently certified final optimality bound proves that true policy regret decreases by at least 58 percent throughout $K$. At the reference state it decreases by at least 75 percent. These are statements about the economic payoff, not about changes in a residual envelope.

The final reference-state optimality bounds lie between $\NineteenBest$ and $\NineteenWorst$, below the unchanged $.01$ target. The corresponding initial-state-uniform bound is $\NineteenK$. Each result concerns newly generated neural time controls, held until the original first exit, with a budget-preserving dependence on initial wealth. The optimal upper comparison retains all originally admissible adapted controls. Neither this policy class nor its initial-state rectangle is identified with the full-state feedback architecture. The latter's best retained complete-domain certificate remains approximately $7.27832$. The original objective of an accurate solution over the entire state--time domain is preserved; a localized certificate is not presented as its completion.

Continuation values are the second substantive issue. The full-state feedback experiments learn critics whose upper-wealth traces remain close to immediate settlement, even though continuation is valuable when that face is inaccessible to the candidate. We evaluate the new policies' one-sided continuation values independently. Their certified excess over immediate settlement exceeds $\NineteenTrace$ uniformly on the tested preference segment. These evaluations provide policy-induced targets for a continuation trace without using the unknown optimal value as a label. A separate, complete-cover intervention raises the old critic's trace while holding its actor fixed. Its global residual bound worsens. The distinction is informative: removing a necessary trace obstruction is not sufficient to construct a globally accurate Bellman upper witness.

The paper also quantifies the contribution of learned initialization in the nonlinear multi-product control problem. On a new, predeclared finite query distribution, neural, zero-start, and accelerated starts are compared at the same directed loss tolerances, rather than after an arbitrary common number of corrections. The learned map reduces gradient work in several regimes, particularly at 8 and 32 state coordinates. At 128 coordinates its advantage is confined to a coarser target; acceleration wins at the strictest target. A learned-gradient-dependent theorem explains the relevant correction budget. The existing cube-uniform guarantee remains explicitly initialization-agnostic, and the new finite-query results are not represented as a cover of that cube.

Finally, economic interpretation is attached to the newly certified neural policies themselves. A wealth increase of $.0092$, financed through an explicit feasible change in consumption, compensates their uniform bound on $K$. It is 0.736 percent of reference wealth. This result is distinct from the retained compensation result for the non-neural price library and does not normalize the much larger full-domain feedback certificate.

\subsection{Contribution and relation to the literature}
Stopped verification, Bellman contraction, strong-convexity bounds, and interval ordering are classical ingredients. \citet{Judd1998} and \citet{KushnerDupuis2001} develop numerical dynamic programming and controlled Markov-chain approximation. \citet{BarlesSouganidis1991} establishes a monotone-approximation framework, and \citet{BrummScheidegger2017} develops adaptive sparse-grid methods for computational economics. Information and control relaxations provide dual comparisons \citep{BrownSmithSun2010,BrownSmith2011}. Neural dynamic-control and differential-equation methods include \citet{HanE2016,HanJentzenE2018}; the retained external comparison uses \citet{CaiFangZhou2025}. Directed arithmetic uses MPFR \citep{FousseEtAl2007}.

Our contribution is the connection of an implemented neural policy-value gate to a coupled-payoff continuum certificate, constructive continuation-trace evaluation, and matched-tolerance accounting for learned representations. The five new actors are not inherited deterministic solutions relabeled as neural networks. Conversely, the faster classical transcription is not hidden, and neither interval ordering nor strong convexity alone is claimed as a new general convergence theorem. The complete-cover objective, accessibility theorem, state--price continuation, dual construction, external holdouts, and adverse historical results are retained. Their roles are separated from the new policy-sensitive results throughout.
